\clearpage
\SourceLeaf{454}%
\pdfbookmark[0]{Abrisse}{abrisse}
\fancyhead[LO]{ABRISSE}
\fancyhead[RE]{ABRISSE}
\begingroup
\fontsize{10.5}{13}\selectfont
\setlength{\tabcolsep}{3pt}
\setlength{\LTpre}{2pt}
\setlength{\LTpost}{2pt}
\begin{longtable}{@{}rrrp{.61\linewidth}@{}}
\multicolumn{4}{@{}p{\linewidth}@{}}{\centering\Large ABRISSE\par\medskip\normalsize DER AUF DEN VERSCHIEDENEN STATIONEN DER GRADMESSUNG 1821. 1822. 1823 UND DEREN FORTSETZUNG BIS JEVER 1824. 1825 FESTGELEGTEN RICHTUNGEN.\par}\\*[18pt]
\multicolumn{4}{c}{\strut STERNWARTE}\\*[5pt]
\multicolumn{4}{c}{\(-5.242\)\qquad\(+0.005\) Theodolithplatz 1821}\\*[5pt]
\multicolumn{4}{c}{\(-5.507\)\qquad\(0\) Theodolithplatz 1823}\\*[5pt]
\multicolumn{4}{@{}p{\linewidth}@{}}{Die Richtungen sind alle auf den Platz von 1823 reducirt.}\\*[5pt]
\(0^{0}\)&\(0'\)&\(2''614\)&\raggedright Südliches Meridianzeichen\tabularnewline
\(10\)&\(12\)&\(42.475\)&\raggedright Meisner Heliotrop\tabularnewline
\(64\)&\(1\)&\(18.020\)&\raggedright Hohehagen (Platz von 1823)\tabularnewline
\(180\)&\(0\)&\(0.000\)&\raggedright Nördliches Meridianzeichen\tabularnewline
\multicolumn{4}{c}{\strut NÖRDLICHES MERIDIANZEICHEN}\\*[5pt]
\multicolumn{4}{c}{\(-5019.756\)\qquad\(-0.133\)}\\*[5pt]
\(0^{0}\)&\(0'\)&\(5''772\)&\raggedright \textit{Sternwarte, Meridianspalt}\tabularnewline
\(0\)&\(23\)&\(54.606\)&\raggedright Hanstein\tabularnewline
\(1\)&\(38\)&\(36.606\)&\raggedright Göttingen, Albani\tabularnewline
\(11\)&\(9\)&\(10.606\)&\raggedright Göttingen, Mariae\tabularnewline
\(13\)&\(9\)&\(0.606\)&\raggedright Weende\tabularnewline
\(18\)&\(9\)&\(22.606\)&\raggedright Klein Schneen\tabularnewline
\(18\)&\(21\)&\(36.606\)&\raggedright Siboldshausen\tabularnewline
\(19\)&\(9\)&\(56.606\)&\raggedright Backhaus Pavillon\tabularnewline
\(21\)&\(20\)&\(26.606\)&\raggedright Rosdorf\tabularnewline
\(27\)&\(11\)&\(1.606\)&\raggedright Volkerode\tabularnewline
\(29\)&\(45\)&\(32.372\)&\raggedright Mengershausen\tabularnewline
\(35\)&\(13\)&\(4.606\)&\raggedright Baum bei Mengershausen\tabularnewline
\(38\)&\(12\)&\(9.606\)&\raggedright Gronde\tabularnewline
\(43\)&\(23\)&\(33.272\)&\raggedright Baum\tabularnewline
\(48\)&\(6\)&\(29.272\)&\raggedright Baum an der Mündner Chaussée\tabularnewline
\(48\)&\(19\)&\(41.527\)&\raggedright \textit{Hohehagen Postament} (1821)\tabularnewline
\multicolumn{3}{r}{\begin{tabular}{@{}rrr|@{}}\(51\)&\(41\)&\(52.606\)\\
\(51\)&\(42\)&\(54.606\)\end{tabular}}&\raggedright Hetgershausen, Kanten des Thurms\tabularnewline
\(64\)&\(33\)&\(26.606\)&\raggedright Elliehausen\tabularnewline
\(101\)&\(41\)&\(24.606\)&\raggedright Lenglern\tabularnewline
\(138\)&\(22\)&\(10.606\)&\raggedright Bovenden\tabularnewline
\(145\)&\(51\)&\(24.606\)&\raggedright Hevensen\tabularnewline
\(148\)&\(22\)&\(51.606\)&\raggedright Wolbrechtshausen\tabularnewline
\(150^{0}\)&\(13'\)&\(56''606\)&\raggedright Parensen\tabularnewline
\(150\)&\(22\)&\(31.606\)&\raggedright Baum auf der Weper\tabularnewline
\(158\)&\(10\)&\(54.606\)&\raggedright Häuschen oberhalb Bovenden\tabularnewline
\(160\)&\(13\)&\(18.606\)&\raggedright Moringen\tabularnewline
\(167\)&\(26\)&\(38.606\)&\raggedright Grossenrode\tabularnewline
\(167\)&\(57\)&\(11.021\)&\raggedright \textit{Hils, Postament}\tabularnewline
\multicolumn{3}{r}{\begin{tabular}{@{}rrr|@{}}\(358\)&\(30\)&\(30.606\)\\
\(358\)&\(32\)&\(36.606\)\\
\(358\)&\(42\)&\(5.606\)\end{tabular}}&\raggedright Kanten des Thibautschen Gartenhauses\tabularnewline
\multicolumn{4}{c}{\strut HOHEHAGEN}\\*[5pt]
\multicolumn{4}{c}{\(+6059.889\)\qquad\(+12447.734\) Hauptplatz von 1821 (1)}\\*[5pt]
\multicolumn{4}{c}{\(+6059.493\)\qquad\(+12448.193\) Nebenplatz von 1822 (2)}\\*[5pt]
\multicolumn{4}{c}{\(+6059.878\)\qquad\(+12447.746\) Platz von 1823 (3)}\\*[5pt]
\multicolumn{4}{@{}p{\linewidth}@{}}{Die beigefügten Zahlen (1), (2), (3) bezeichnen die Standpunkte, von wo aus die Schnitte gemacht sind, die mit Cursivbuchstaben bezeichneten Richtungen sind am Platz (3) gemachte oder darauf reducirte Schnitte.}\\*[5pt]
\(3^{0}\)&\(47'\)&\(52''920\)&\raggedright Meensen (3)\tabularnewline
\(55\)&\(59\)&\(40.490\)&\raggedright \textit{Hercules}\tabularnewline
\(64\)&\(0\)&\(39.064\)&\raggedright Burghasungen (1)\tabularnewline
\(41\)&\(43\)&\(53.800\)&\raggedright Landwehrhagen (1)\tabularnewline
\(41\)&\(57\)&\(7.800\)&\raggedright Lutternberg (1)\tabularnewline
\(165\)&\(22\)&\(49.800\)&\raggedright Wolfstrang (1)\tabularnewline
\(185\)&\(48\)&\(16.262\)&\raggedright \textit{Hils}\tabularnewline
\(186\)&\(37\)&\(13.800\)&\raggedright Hube, Durchschnitt (1)\tabularnewline
\(193\)&\(32\)&\(13.161\)&\raggedright Ochsenberg (3)\tabularnewline
\(197\)&\(34\)&\(49.298\)&\raggedright \textit{Beinberg} (1)\tabularnewline
\(211\)&\(2\)&\(6.155\)&\raggedright Echte (?) (2)\tabularnewline
\(212\)&\(19\)&\(44.155\)&\raggedright Nordheim, kleiner Thurm (2)\tabularnewline
\(212\)&\(31\)&\(26.155\)&\raggedright Nordheim, Rathhaus (2)\tabularnewline
\(212\)&\(50\)&\(7.155\)&\raggedright Nordheim, Kirchthurm (2)\tabularnewline
\(225\)&\(33\)&\(53.612\)&\raggedright Plesse dünner Thurm (1 u. 3)\tabularnewline
\multicolumn{3}{r}{\begin{tabular}{@{}rrr|@{}}\(225\)&\(37\)&\(14.064\)\\
\(225\)&\(38\)&\(26.064\)\end{tabular}}&\raggedright Hetgershausen, Kanten des Thurms (1)\tabularnewline
\(225\)&\(38\)&\(14.161\)&\raggedright Hetgershausen, Fahnenstange (3)\tabularnewline
\(226\)&\(40\)&\(18.392\)&\raggedright Windmühle bei Clausthal (1)\tabularnewline
\end{longtable}
\endgroup
